\documentclass[12pt]{article}
\usepackage[T1]{fontenc}
\usepackage{mathpazo}
\usepackage[margin=1in]{geometry}
\usepackage[hidelinks]{hyperref}
\usepackage{setspace}
\pagestyle{empty}
\begin{document}
\begin{center}
\renewcommand{\thefootnote}{\fnsymbol{footnote}}
{\Large\bfseries Housing First or Treatment First? Ask the Kids, Not the Parents}\footnote{I thank Ay\c{s}e \.{I}mrohoro\u{g}lu, Kai Zhao, Joseph Steinberg, Efraim Benmelech, Jonathan Alschech, Ryan Quandt, Christopher Sharp and the seminar participants at the University of Toronto for their helpful comments. The data and replication code are available on GitHub at \url{https://github.com/srotberg/housing_first_or_treatment_first_ask_the_kids_not_the_parents}.}\\[1.4em]
{\large Shahar Rotberg}\footnote{Canada Mortgage and Housing Corporation (CMHC). Email: \url{shahar.rotberg@gmail.com}. Website: \url{www.shaharrotberg.com}. This work was conducted on personal equipment and time, was not funded by any person or entity, and the author was not influenced by anyone at any point in this research. The author owns the code; he wrote the baseline code of the quantitative model. Claude (Anthropic) was used to write the code of the stylized model and extensions of the quantitative model, to debug the code, and to run the simulations. Every number in the paper is generated from that code, the two models were checked against each other on every mechanism they share, with the discrepancies found and their resolution logged in the replication package, and the code is public. The opinions expressed herein do not necessarily represent the views of CMHC.}\\[0.6em]
Canada Mortgage and Housing Corporation\\
\href{mailto:shahar.rotberg@gmail.com}{shahar.rotberg@gmail.com}\\[1em]
November 2026
\end{center}
\renewcommand{\thefootnote}{\arabic{footnote}}
\setcounter{footnote}{0}

\vspace{1.5em}
\begin{center}
\begin{minipage}{0.88\textwidth}
\begin{center}\textbf{Abstract}\end{center}
\setstretch{1.15}
\noindent Should housing assistance require sobriety? A general-equilibrium model of homelessness, drug use, drug prices and mortality is calibrated to US data. Use concentrates among low-income households because a substance use disorder's mortality risk, unlike its unemployment risk, is uninsurable. Drug use harms users' children, cutting their consumption and, through a parent's disorder, their adult earnings. In the long run, sobriety conditions reduce these harms and, in aggregate, also childhood homelessness. Consequently, a newborn's welfare rises by $2.57$ percent of lifetime consumption, of which three quarters is an externality its parent does not count, and in general equilibrium parents gain too. Even along the transition, homelessness and childhood homelessness do not rise, and newborns still gain. The model reproduces micro evidence of randomised controlled trials, and users with a disorder lose, but macro effects on other households and children overturn the verdict. Moral hazard alone does not justify conditionality.

\vspace{1em}
\noindent\textit{Keywords:} homelessness; housing vouchers; substance use disorder; intergenerational spillovers; externalities.

\vspace{0.4em}
\noindent\textit{JEL classification:} I18; I38; H23; R21; J62.
\end{minipage}
\end{center}
\end{document}
